\section{Self-improvement behind one interface}
\label{sec:methods}

PyOD put outlier detectors behind one \texttt{fit} / \texttt{decision\_function} interface
\citep{zhao2019pyod}. That interface later let ADBench compare 30 detectors on equal terms
\citep{han2022adbench}. \lib{} does the same for self-improvement. Every method is a class
with one constructor and three calls:
\begin{lstlisting}[language=Python]
m = rsi.Reflexion(library=True,
                  cards=True)
m.fit(generate, tasks.dev)
ep = m.solve(generate, task, seed=11)
ep.accepted, ep.code, ep.usage
m.state()
\end{lstlisting}
\texttt{generate} is the user's own model. \texttt{fit()} does nothing for a method that does
not learn. \texttt{state()} returns what the method learned,
which is hashed into every episode record. \texttt{rsi.benchmark()} fits each method once on
the development tasks, runs it on every held-out task and seed, and grades the results.

\paragraph{What a method sees (Figure~\ref{fig:overview}b).}
\texttt{solve()} receives the task text and a judge, on which it can only call
\texttt{check(code)}. A check returns whether the program ran, its error, the result of each
check, and the causes of the failed ones. It never returns the grade: whether the final
program matches an independent reference. The benchmark runner computes the grade after
\texttt{solve()} returns and never passes it back (\S\ref{sec:grader}). Learning methods may
use development grades to keep or reject a candidate, but no text a model reads contains a
grade.

\paragraph{Equal budgets.}
Every model call goes through a meter, which refuses a call beyond the budget and records
which limit stopped the episode. Each episode gets:
\begin{itemize}
\item \BudgetAttempts{} program executions;
\item \BudgetCalls{} model calls, counting helper calls such as reflections and reviews;
\item \BudgetTokens{} generated tokens.
\end{itemize}
Learning methods also get a development budget of \FitEpisodes{} episodes, \FitProposals{}
proposals, \FitCalls{} calls and \FitTokens{} tokens. Development episodes count against
both budgets. Token counts come from the backend when it reports them; otherwise they are
estimated as one token per four characters.

\begin{table*}[t]
\centering
\small
\setlength{\tabcolsep}{4pt}
\begin{tabular}{@{}lp{5.3cm}p{3.9cm}p{2.8cm}@{}}
\toprule
Method & At test time & Learned on development tasks & In the spirit of \\
\midrule
\texttt{NoImprove} & one program, no feedback & -- & baseline \\
\texttt{SelfDebug} & retry with the error only; stop once the program runs & -- &
\citet{chen2024selfdebug} \\
\texttt{CheckFeedback} & retry with the failed checks, their causes and cards & -- &
verifier feedback \\
\texttt{Reflexion} & as \texttt{CheckFeedback}, plus a written reflection before each retry &
-- & \citet{shinn2023reflexion} \\
\texttt{PlaybookLearning} & as \texttt{CheckFeedback}, with a frozen playbook in the prompt &
playbook entries, through the gated loop & \citet{zhao2024expel,suzgun2026dynamic} \\
\texttt{Team} & a writer and one reviewer per rule family, from the same model & -- &
multi-role review \\
\texttt{Evolve} & as \texttt{CheckFeedback}, with an evolved prompt addendum & the addendum,
over \EvolveGenerations{} generations on \EvolveMinibatch{} tasks & \citet{agrawal2026gepa} \\
\bottomrule
\end{tabular}
\caption{The \NMethods{} self-improvement methods. All share the prompt, the judge and the
budget; each keeps the mechanism it is known for. They are not reproductions of the original
systems.}
\label{tab:methods}
\end{table*}

\paragraph{Methods.}
Table~\ref{tab:methods} lists the \NMethods{} methods.
\begin{itemize}
\item \texttt{SelfDebug} returns only the error of a program that did not run, and stops as
soon as a program runs.
\item \texttt{CheckFeedback} also returns the failed checks.
\item \texttt{Reflexion} first asks the model for a reflection of at most five sentences,
without code, and keeps the reflections for the rest of the episode.
\item In \texttt{Team}, a reviewer is asked only when a check of its own family failed. It sees
the program, the failed checks, their causes and its family's cards, and answers with a
diagnosis rather than code. A program that does not run goes to an interface reviewer.
\item \texttt{Evolve} rewrites a prompt addendum of at most 1,200 characters from the visible
failures of the current best addendum. It keeps a candidate only if it scores strictly higher
on the same development tasks: first by the number of valid programs, then by the share of
passed checks.
\end{itemize}

\paragraph{Gated playbook loop.}
The playbook method is built on the library's \texttt{RSILoop}. A playbook is a list of short
entries on three editable surfaces: tool usage, retrieval and experience. The proposer sees
development traces only. For each failed attempt it sees the program, the failed checks and
their causes, but never the grade. One proposal changes one entry. The candidate and the
current playbook are both evaluated on two disjoint halves of the development tasks, split by
algorithm. The candidate is accepted only if all of the following hold:
\begin{itemize}
\item its score does not drop on either half;
\item its score rises on at least one half;
\item no task that passed every check before fails one under the candidate.
\end{itemize}
The last condition is the \emph{no-regression} rule. The stricter \emph{all-pass} rule vetoes
any candidate while a single task still fails, which blocks learning for a weak model.
Whitespace-only edits and candidates that were already tried are rejected. Every proposal,
evaluation and decision is appended to a journal whose receipts carry evidence hashes.
Timestamps may not move backward, and a frozen playbook cannot change.

\subsection{Skill arms}
\label{sec:arms}

Two switches separate the parts of the skills. \texttt{library} adds the interface text and
allows \texttt{import fin\_skills}; \texttt{cards} adds the task's rule cards to the prompt
and to the check feedback. Together they define \NArms{} arms:

\begin{center}
\small
\begin{tabular}{@{}lcc@{}}
\toprule
Arm & Rule cards & Library \\
\midrule
\texttt{none} & -- & -- \\
\texttt{cards} & \checkmark & -- \\
\texttt{library} & -- & \checkmark \\
\texttt{cards+library} & \checkmark & \checkmark \\
\bottomrule
\end{tabular}
\end{center}

Every arm keeps the task definition: the rule, the timing convention and the form of the split
factors. These fix the correct answer and are not skills. Every arm also sees the failed checks
and their causes. Within one arm, \texttt{NoImprove} versus \texttt{CheckFeedback} measures the
checks, and \texttt{CheckFeedback} versus \texttt{PlaybookLearning} measures self-improvement.
For one method, comparing arms measures the skills.

Each arm sees only its own material. Without the library, the worker refuses
\texttt{import fin\_skills}, and cards are rendered without library calls. The development
traces given to a proposer follow the same arm.
